\documentclass[11pt]{article}
\usepackage[margin=1in]{geometry}
\usepackage{amsmath,amssymb,amsthm}
\usepackage{hyperref}

\newtheorem{theorem}{Theorem}
\newtheorem{lemma}[theorem]{Lemma}
\newtheorem{proposition}[theorem]{Proposition}
\newtheorem{corollary}[theorem]{Corollary}
\theoremstyle{definition}
\newtheorem{definition}[theorem]{Definition}
\theoremstyle{remark}
\newtheorem{remark}[theorem]{Remark}

\newcommand{\tors}{\operatorname{tors}}
\newcommand{\add}{\operatorname{add}}
\newcommand{\Filt}{\operatorname{Filt}}
\newcommand{\Hom}{\operatorname{Hom}}
\newcommand{\Ext}{\operatorname{Ext}}
\newcommand{\md}{\operatorname{mod}}

\title{Disproof of Conjecture 00000009898:\\
the torsion lattice of $K(1\to 2)$ is the pentagon}
\author{jilint777}
\date{2026-10-04}

\begin{document}
\sloppy
\maketitle

\begin{abstract}
Conjecture 00000009898 asserts that the lattice $\tors A$ of torsion classes of a
finite-dimensional algebra $A$ is distributive \emph{if and only if} the oriented path partial
order of the quiver of $A$ is a forest order, and that the minimal counterexample to
distributivity is the three-vertex oriented cycle. Both claims are false. Over \emph{every} field
$K$, the path algebra $A=KA_2$ of the quiver $1\to 2$ has a 2-element chain as its path order,
which is a forest order under every reading. Yet $\tors A$ is the pentagon $N_5$, which is not
distributive, and $A_2$ has only two vertices. The refutation is fully formalised in Lean~4
(core library only, for an arbitrary field). It defines representations, morphisms, quotients,
short exact sequences, torsion classes and the lattice operations from scratch. An independent
Python script recomputes the whole lattice by brute force over $\mathbb F_2$ and $\mathbb F_3$.
We also prove the correct criterion: $\tors A$ is distributive iff the quiver of $A$ has no arrow
between distinct vertices. So the ``only if'' direction of the conjecture holds, and the ``if''
direction fails for every quiver with such an arrow.
\end{abstract}

\section{The claim}

The conjecture (English version): \emph{``The torsion class lattice of a finite-dimensional
algebra is distributive if and only if the oriented path partial order of the quiver is a forest
order; the minimal counterexample to distributivity is the three-vertex oriented cycle, and the
length of distributive approximations of the lattice is bounded.''}

The statement is a conjunction of three clauses:
\begin{itemize}
\item[(C1)] for every finite-dimensional algebra $A$: $\tors A$ is distributive $\iff$ the path
order of the quiver $Q_A$ is a forest order;
\item[(C2)] the minimal counterexample to distributivity is the oriented $3$-cycle. In particular,
every algebra whose quiver has fewer than $3$ vertices has a distributive torsion lattice;
\item[(C3)] ``the length of distributive approximations is bounded''. This is not defined in the
statement, and we do not use it.
\end{itemize}
We refute the direction ``$\Leftarrow$'' of (C1), hence (C1), and also (C2). One false clause
already makes the conjunction false.

\section{Set-up}

Let $K$ be an arbitrary field and $A=KA_2$ the path algebra of $A_2\colon 1\xrightarrow{\ a\ }2$.
It has dimension $3$, with basis $e_1,e_2,a$. As usual, finite-dimensional $A$-modules are the
same as finite-dimensional representations $M=(V_1\xrightarrow{f}V_2)$: finite-dimensional
$K$-spaces $V_1,V_2$ with a $K$-linear map $f$. A morphism $(g_1,g_2)\colon M\to N$ is a pair of
linear maps with $f_N g_1=g_2 f_M$. Kernels, images and cokernels are computed vertexwise, so
$M\to N$ is an epimorphism (resp.\ monomorphism) iff $g_1,g_2$ are surjective (resp.\ injective),
and a sequence $0\to L\to M\to N\to 0$ is exact iff it is exact at both vertices.

The indecomposable representations are
\[
S_1=(K\to 0),\qquad S_2=(0\to K),\qquad P_1=(K\xrightarrow{\ \mathrm{id}\ }K).
\]
Here $S_1,S_2$ are the simples, $P_2=S_2$, and $P_1=I_2$ is projective-injective.

\begin{definition}
A \emph{torsion class} is a class $\mathcal T$ of modules that contains $0$ and is closed under
quotients (images of epimorphisms) and extensions. Closure under isomorphisms is a special case
of closure under quotients. $\tors A$ is the set of torsion classes ordered by inclusion. It is a
complete lattice: the meet is the intersection, and the join $\mathcal T\vee\mathcal U$ is the
intersection of all torsion classes containing $\mathcal T\cup\mathcal U$.
\end{definition}

If one does not require $0\in\mathcal T$, the only new torsion class is $\emptyset$, since a nonempty class closed under
quotients contains $0$ (as a quotient $M\twoheadrightarrow 0$). All classes used
below contain $0$, so the argument is unchanged.

\paragraph{The path order.} For a quiver $Q$ write $i\le j$ if there is an oriented path (of
length $\ge0$) from $i$ to $j$. For $A_2$ this is the chain $1<2$. Some authors orient the other
way, which gives $2<1$, again a chain. A \emph{forest order} is defined in the literature in
several ways: every principal down-set is a chain (a rooted forest); every principal up-set is a
chain; or the Hasse diagram has no cycle. A 2-element chain satisfies all of these, and it is a
partial order. The quiver of $A=KA_2$ (its Gabriel quiver) is $A_2$ itself, since $A$ is basic
with $\operatorname{rad}A/\operatorname{rad}^2A=Ka$. Hence
\begin{center}
\emph{the path order of the quiver of $KA_2$ is a forest order, for every reading.}
\end{center}

\section{The five torsion classes}

Define three classes of representations:
\[
x=\{M: V_2=0\},\qquad y=\{M: f\text{ surjective}\},\qquad z=\{M: V_1=0\}.
\]
Also let $0$ be the class of zero representations, and $\md A$ the class of all of them.

\begin{lemma}\label{lem:tors}
$x$, $y$ and $z$ are torsion classes.
\end{lemma}
\begin{proof}
All three contain $0$.

\emph{$x$.} Let $(g_1,g_2)\colon M\twoheadrightarrow N$ with $M\in x$. Then
$V_2^N=g_2(V_2^M)=g_2(0)=0$. Let $0\to L\xrightarrow{i}M\xrightarrow{p}N\to0$ be exact with
$L,N\in x$, and let $v\in V_2^M$. Then $p_2(v)\in V_2^N=0$, so $v=i_2(u)$ for some
$u\in V_2^L=0$, and $v=0$.

\emph{$z$.} The same argument at vertex $1$.

\emph{$y$, quotients.} Let $w\in V_2^N$. Write $w=g_2(v)$, and $v=f_M(u)$. Then
$w=g_2f_M(u)=f_N(g_1u)$.

\emph{$y$, extensions.} Let $w\in V_2^M$. Since $N\in y$, we have $p_2(w)=f_N(a)$ for some $a$.
Since $p_1$ is onto, $a=p_1(b)$. Then $p_2(w)=f_Np_1(b)=p_2f_M(b)$, so
$w-f_M(b)\in\ker p_2=\operatorname{im} i_2$, say $w-f_M(b)=i_2(c)$. Since $L\in y$, $c=f_L(d)$,
so $i_2(c)=i_2f_L(d)=f_M(i_1d)$. Therefore $w=f_M(b+i_1d)$.
\end{proof}

\begin{lemma}\label{lem:all}
Every torsion class $\mathcal T$ with $S_1,S_2\in\mathcal T$ is $\md A$.
\end{lemma}
\begin{proof}
The sequences $0\to S_i\to S_i^{n+1}\to S_i^n\to0$ (inclusion of the first coordinate,
projection onto the others) are exact. By induction, $S_1^m,S_2^n\in\mathcal T$ for all $m,n$;
the case $n=0$ is $0\in\mathcal T$. For an arbitrary $M=(K^{m}\xrightarrow{f}K^{n})$ the sequence
\[
0\to(0\to K^n)\xrightarrow{(0,\mathrm{id})}M\xrightarrow{(\mathrm{id},0)}(K^m\to0)\to0
\]
is exact. Its outer terms are $S_2^n$ and $S_1^m$, so $M\in\mathcal T$.
\end{proof}

\begin{proposition}\label{prop:N5}
$\tors(KA_2)=\{0,\ x,\ y,\ z,\ \md A\}$. These classes form the pentagon $N_5$:
\[
0\subsetneq x\subsetneq y\subsetneq \md A,\qquad 0\subsetneq z\subsetneq\md A,
\]
with $z$ incomparable to $x$ and to $y$. Also $x=\add S_1$, $y=\add\{S_1,P_1\}$ and
$z=\add S_2$.
\end{proposition}
\begin{proof}
If $f\colon V_1\to V_2$ has rank $r$, choose a complement $C$ of $\ker f$ and a complement $D$
of $\operatorname{im}f$. In adapted bases, $M\cong P_1^r\oplus S_1^{\dim V_1-r}\oplus
S_2^{\dim V_2-r}$. A torsion class is closed under direct sums (split extensions) and under
direct summands (quotients). So it is $\add$ of the set of indecomposables it contains. That set
$I\subseteq\{S_1,S_2,P_1\}$ satisfies $P_1\in I\Rightarrow S_1\in I$, because
$P_1\twoheadrightarrow S_1$. It also satisfies $S_1,S_2\in I\Rightarrow P_1\in I$, by
Lemma~\ref{lem:all} or by the extension $0\to S_2\to P_1\to S_1\to0$. This leaves the
possibilities $\emptyset$, $\{S_1\}$, $\{S_2\}$, $\{S_1,P_1\}$ and $\{S_1,S_2,P_1\}$. They give
$0$, $x$, $z$, $y$ and $\md A$, which are torsion classes by Lemma~\ref{lem:tors}. The
inclusions are witnessed by $S_1\in x\setminus 0$, $P_1\in y\setminus x$,
$S_2\in z\setminus y$ and $S_1\in x\setminus z$.
\end{proof}

Only Lemmas~\ref{lem:tors} and~\ref{lem:all} are needed for the disproof. The classification in
Proposition~\ref{prop:N5} shows that the counterexample is the whole lattice, not just a
sublattice.

\section{Failure of distributivity}

\begin{theorem}\label{thm:main}
For every field $K$, $\tors(KA_2)$ is not distributive:
\[
y\wedge(x\vee z)=y\neq x=(y\wedge x)\vee(y\wedge z).
\]
The dual law fails as well: $x\vee(y\wedge z)=x\neq y=(x\vee y)\wedge(x\vee z)$.
\end{theorem}
\begin{proof}
By Lemma~\ref{lem:all}, $x\vee z=\md A$, because $S_1\in x$ and $S_2\in z$. So
$y\wedge(x\vee z)=y$. Next, $y\wedge x=x$ since $x\subseteq y$ (if $V_2=0$, $f$ is onto). Also
$y\wedge z=0$: if $V_1=0$ and $f$ is onto, then $V_2=f(0)=0$. Since $x$ is a torsion class
containing $x\cup 0$, we get $(y\wedge x)\vee(y\wedge z)=x$. Finally $P_1\in y\setminus x$. For
the dual law, $x\vee y=y$ and $x\vee 0=x$.
\end{proof}

\begin{corollary}
Clause (C1) is false: for $A=KA_2$ the path order is a forest order but $\tors A$ is not
distributive. Clause (C2) is false: $A_2$ has $2<3$ vertices and a non-distributive torsion
lattice. Since a one-vertex algebra $A$ is local, with $\tors A=\{0,\md A\}$ (a $2$-chain), the
true minimum number of vertices is $2$.
\end{corollary}

\section{Readings covered}

\begin{itemize}
\item \textbf{Field.} Everything holds over every field. The Lean proof is for an arbitrary
field, with instances $\mathbb F_2$ and $\mathbb F_3$.
\item \textbf{Forest order.} The path order of $A_2$ is a 2-chain: a partial order whose down-sets
and up-sets are chains and whose Hasse diagram is a single edge. It is a forest under every
definition and for either orientation convention (left or right modules, i.e.\ $1<2$ or
$2<1$). Lean uses the strongest version (antisymmetric, and both down-sets and up-sets are chains).
\item \textbf{Distributivity.} Both distributive laws fail, so any definition of a distributive
lattice is covered.
\item \textbf{Torsion classes.} The result is the same whether or not $\emptyset$ is admitted, and
closure under isomorphism is automatic.
\item \textbf{Minimality (C2).} Measured by vertices (or arrows), $A_2$ is smaller than the
3-cycle. If ``minimal counterexample'' means a minimal quiver whose path order is \emph{not} a
forest, the 2-cycle $1\rightleftarrows2$ is smaller still. With relations
$\operatorname{rad}^2=0$ this is the preprojective algebra of $A_2$, whose torsion lattice is the
weak order of $S_3$, a hexagon.
\item \textbf{(C3)} is undefined. The disproof does not depend on it.
\end{itemize}

\section{Remark: the correct criterion}

The following is not needed for the disproof. It shows that only the ``if'' half of (C1) fails.

\begin{proposition}
Let $A$ be a finite-dimensional basic algebra over a field. Then $\tors A$ is distributive if and
only if the quiver of $A$ has no arrow between two distinct vertices. In that case the path order
is an antichain, hence a forest order.
\end{proposition}
\begin{proof}
Suppose there is an arrow between distinct vertices $i\ne j$. Then
$\Ext^1(S_i,S_j)\ne0$, possibly after exchanging $i$ and $j$ according to the convention, so
there is a non-split extension $0\to S_j\to M\to S_i\to0$. For a class $\mathcal C$ let
$T(\mathcal C)$ be the smallest torsion class containing it. The class
$\{X:\Hom(X,S_j)=0\}$ is closed under quotients and extensions (left exactness of
$\Hom(-,S_j)$), and it contains $M$: a nonzero map $M\to S_j$ would split the sequence. Hence
$S_j\notin T(M)$, and since every nonzero $X\in T(S_j)\subseteq\Filt(S_j)$ has
$\Hom(X,S_j)\ne0$, we get $T(M)\wedge T(S_j)=0$. Moreover $S_i\in T(M)$, and $M\in T(S_i)\vee
T(S_j)$. But $M\notin T(S_i)\subseteq\Filt(S_i)$. Thus
$T(M)\wedge(T(S_i)\vee T(S_j))=T(M)\ne T(S_i)=(T(M)\wedge T(S_i))\vee(T(M)\wedge T(S_j))$.

Conversely, if no arrow joins distinct vertices, then $\Ext^1(S_i,S_j)=0$ for $i\ne j$. The
blocks of $A$ correspond to the connected components of its quiver, so $A=\prod_iA_i$ with each
$A_i$ local. A nonzero torsion class of a local algebra contains the unique simple module, hence
all modules. So $\tors A\cong\prod_i\{0<1\}$ is Boolean, hence distributive.
\end{proof}

So the ``only if'' direction of (C1) is true. The ``if'' direction fails for every quiver with an
arrow between distinct vertices, and $A_2$ is the smallest example.

\section{Machine verification}

\paragraph{Lean 4 (\texttt{lean4/Main.lean}, core only, v4.19.0).} Everything is defined from
scratch:
\begin{itemize}
\item a field structure \texttt{Fld} (all field axioms), coordinate spaces $K^n$ and linear maps;
\item representations \texttt{Rep K} of $A_2$ (dimensions $d_1,d_2$ arbitrary, $f$ any linear
map), morphisms \texttt{Hom}, epimorphisms, isomorphisms, short exact sequences;
\item \texttt{IsTors} (contains $0$, closed under quotients and extensions), meets and joins, with
proofs that they are torsion classes and are the greatest lower and least upper bounds;
\item the quiver $A_2$, oriented paths, \texttt{ForestOrder}.
\end{itemize}
Main results: \texttt{xC\_isTors}, \texttt{yC\_isTors}, \texttt{zC\_isTors} (Lemma~\ref{lem:tors},
for all dimensions); \texttt{all\_of\_S1\_S2} (Lemma~\ref{lem:all}); \texttt{pentagon} (the $N_5$
relations); \texttt{not\_distributive}, \texttt{not\_codistributive} (Theorem~\ref{thm:main});
\texttt{a2\_forest}; and
\[
\texttt{conjecture\_00000009898\_false (K) [Fld K] : }\neg\,\texttt{IfClause K},
\]
where \texttt{IfClause K := ForestOrder A2Q → Distributive K}. Also \texttt{iffClause\_false} and
\texttt{minClause\_false}.

Non-vacuity: \texttt{five\_tors} shows that $0,x,y,z,\md A$ are torsion classes and that
$\{f\text{ injective}\}=\add\{S_2,P_1\}$ is not. Field instances are provided for $\mathbb F_2$
and $\mathbb F_3$. \texttt{c3\_facts} records that the 3-cycle has 3 vertices and a
non-antisymmetric path relation. The axioms used are only \texttt{propext}, \texttt{Quot.sound}
and \texttt{Classical.choice}.

\emph{Limitations.} Representations use coordinate spaces $K^{d}$. Every finite-dimensional
space is isomorphic to one, and torsion classes are closed under isomorphism, so this loses
nothing. The equivalence between $KA_2$-modules and representations of $A_2$, the classification
in Proposition~\ref{prop:N5}, and the general criterion are proved only in this report.

\paragraph{Python (\texttt{verify.py}, standard library).} A different method: over $\mathbb F_2$
and $\mathbb F_3$ it enumerates all representations with $d_1,d_2\le2$ and computes isomorphism
classes by brute force under $GL\times GL$ (14 classes). It finds the quotient relation by
searching for surjective morphisms, and the extension triples from kernels of surjective
morphisms. It then tests all $2^{14}$ subsets for being torsion classes. It finds exactly five,
equal to $0,x,y,z,\md A$, forming $N_5$. Distributivity fails, and exhaustively 2 of the 125
triples violate it. It also checks the path orders of $A_2$ and of the 3-cycle.

\end{document}
